% Levenberg-Marquardt correction: the shift delta I moves every eigenvalue by delta.
%
%   figures/tikz/qn-eigen-shift.tex
%       -> media/figures/qn-eigen-shift.{png,svg}
%
% Lecture 16 (Quasi-Newton Methods), Section III.B, before the activity on choosing delta.
%
% GEOMETRY (1 unit = 0.85 cm). Eigenvalues of the Hessian: lambda = (-2, 0.5, 3).
% Floor epsilon = 1. delta = -min_j(lambda_j - epsilon, 0) = -min(-3, -0.5, 2, 0) = 3.
% Shifted eigenvalues lambda + delta = (1, 3.5, 6): the smallest lands exactly on
% epsilon. Arithmetic by hand, 2026-10-07.
%
% GREYSCALE: Hessian eigenvalues are open circles, shifted ones filled; the shift is an
% arrow; the epsilon floor is a dashed line; zero is a solid tick. Labelled throughout.
\definecolor{okabeblue}{HTML}{0072B2}

\begin{tikzpicture}[
    x=0.85cm, y=1cm,
    >={Latex[length=2.0mm]},
    font=\small,
    lab/.style={font=\scriptsize, inner sep=1.5pt},
  ]
  % two number lines
  \draw[black!55, ->] (-2.8,1.2) -- (6.8,1.2) node[right, lab] {$\lambda_j$ of $\nabla^2 f$};
  \draw[black!55, ->] (-2.8,0) -- (6.8,0) node[right, lab] {$\lambda_j + \delta$};
  % zero and epsilon
  \foreach \yy in {0,1.2} { \draw[black!55] (0,\yy-0.1) -- (0,\yy+0.1); }
  \node[lab, below] at (0,-0.08) {$0$};
  \draw[okabeblue, dashed, thick] (1,-0.45) -- (1,1.55);
  \node[lab, okabeblue, above] at (1,1.55) {$\epsilon$};
  % Hessian eigenvalues (open) and shifted eigenvalues (filled)
  \foreach \l/\s in {-2/1, 0.5/3.5, 3/6} {
    \draw[->, thick] (\l,1.08) -- (\s,0.12);
    \filldraw[fill=white] (\l,1.2) circle (2.2pt);
    \filldraw[black] (\s,0) circle (2.2pt);
  }
  \node[lab, above] at (-2,1.28) {$-2$};
  \node[lab, above] at (0.5,1.28) {$0.5$};
  \node[lab, above] at (3,1.28) {$3$};
  \node[lab, below] at (3.5,-0.08) {$3.5$};
  \node[lab, below] at (6,-0.08) {$6$};
  \node[lab, anchor=west] at (5.6,0.75) {shift $\delta = 3$};
\end{tikzpicture}
